\documentclass[a4paper]{article}
\usepackage{amsfonts}
\usepackage{amsmath}
\usepackage{amssymb}
\usepackage{graphicx}      

\newcommand{\dint}{\displaystyle\int}
\newcommand{\dlim}{\lim\limits}

\begin{document}
  
\noindent \large Auteur: Thijs Van Schel \\
\noindent \large Studierichting: Informatica\\
\noindent \large Oefeningenreeks 7E nr. 2.1\\

\medskip

\normalsize

$f(x,y) = (x-3)^6 + (2x+y)^8$ \\

$
\left\{ 
\begin{array}{l} 
\dfrac{\partial f}{\partial x} = 6(x-3)^5 + 16(2x+y)^7 = 0 \\ 
\dfrac{\partial f}{\partial y} = 8(2x+y)^7 = 0 
\end{array} 
\right.
\Rightarrow
\left\{ 
\begin{array}{l} 
6(x-3)^5 + 16(0) = 0 \\ 
2x+y = 0 
\end{array} 
\right.
\Rightarrow
\left\{ 
\begin{array}{l} 
x = 3 \\ 
y = -6 
\end{array} 
\right.
$\\

$\Rightarrow (x,y) \in \{(3, -6)\}$\\

$
\left\{ 
\begin{array}{l} 
\dfrac{\partial^2 f}{\partial x^2} = 30(x-3)^4 + 224(2x+y)^6 \\ 
\dfrac{\partial^2 f}{\partial x \partial y} = \dfrac{\partial^2 f}{\partial y \partial x} = 112(2x+y)^6 \\
\dfrac{\partial^2 f}{\partial y^2} = 56(2x+y)^6
\end{array} 
\right.
\Rightarrow H(x,y) = \left(\begin{array}{c c} 0 & 0 \\ 0 & 0 \end{array}\right) \text{ voor } (3,-6)
$\\

$H(3, -6) = 0 \Rightarrow$ geen uitsluitsel\\

$f(x,y) = (x-3)^6 + (2x+y)^8 \ge 0$ voor alle $(x,y)$\\

$f(3,-6) = 0$ is de kleinst mogelijke waarde\\

$\Rightarrow (3, -6) = $ absoluut minimum\\

\begin{thebibliography}{999}
%\bibitem{a} Referentie
\end{thebibliography}
\end{document}